% Experimental setup: what the experiments of Sections 5 and 6 share. Experiment-specific choices (baselines, triggers,
% the audit's protocol) stay with their experiments.

The experiments of Sections~\ref{sec:XML:Pullback} and~\ref{sec:XML:Pushforward} share an instantiation of $P$, data,
detectors, explainers, and a way of measuring faithfulness. We describe them once here; the baselines, the planted
triggers, and the audit's protocol belong to individual experiments and are described with them.

\subsection{Provenance for EMBER}
\label{sec:XML:SetupEMBER}

\begin{table}[t]
  \centering
  \small
  \begin{tabular}{@{}lrl@{}}
    \toprule
    Feature group & Features & Produced by \\
    \midrule
    Header fields\textsuperscript{a} & 149 & the bytes of the field \\
    Byte histogram & 256 & every byte, in the bin of its value \\
    Byte-entropy histogram & 256 & windowed bytes, binned \\
    Strings & 177 & the strings counted or summarized \\
    Sections & 224 & sections, hashed by name \\
    Imports, exports & 1{,}411 & names of functions, hashed \\
    General & 7 & the whole file, or none \\
    Parser warnings & 88 & none (the null element) \\
    \bottomrule
  \end{tabular}
  \caption{Provenance in the instrumented EMBER v3 extractor. Fields, counts, histograms, and hashed counts are additive
  over their items and get their Shapley provenance (Proposition~\ref{prop:AdditiveShapley}); entropies and averages
  follow the data-flow convention of Section~\ref{sec:XML:Constructing}. \textsuperscript{a}Including data directories,
  the Rich header and Authenticode fields.}
  \label{tab:XML:Conventions}
\end{table}

We instrument the reference implementation of the EMBER v3 extractor~\cite{joyce2025ember}, which computes 2{,}568
features from a PE file. Table~\ref{tab:XML:Conventions} lists, for each group of features, the items it reads. Most
groups are additive over their items and receive their Shapley provenance: a header field is one item, the bytes it
occupies; the byte histogram counts every byte into the bin of its value; and section, import, and export features hash
names into buckets and count them. The byte-entropy histogram and the summary statistics of strings, such as the
average string length, are not additive and follow the data flow: each window or string that a feature is computed
from receives one unit.
Features that no byte produces go to the null element: the parser warnings, whose source the reference extractor does
not report, and, file by file, every feature whose value records an absence, such as a hashed import bucket into which
no imported function falls.

The instrumentation adds 616 lines to the extractor's 1{,}104 and changes none of them. The instrumented extractor
reproduces the original features exactly, except for one hash-seed-dependent quirk of the reference code, and extracts
the features with their provenance in 0.15\,s for the median file (208\,KB). Section~\ref{sec:XML:PullbackFaithfulness}
compares it with the black-box estimate of Section~\ref{sec:XML:BlackBoxP}.

\subsection{Data, Detectors, and Explainers}
\label{sec:XML:SetupData}

We use the PE files of EMBER-2024~\cite{joyce2025ember}, truncated to their first 4{,}000{,}000 bytes (13--16\% of the
files are longer), so that the extractor, $P$, and every detector see the same object. The \emph{explanation set} is
2{,}000 test files, 1{,}000 per label, deduplicated with TLSH and with at most 10 files per malware family; unless stated
otherwise, experiments use a subset of 100 of them (50 per label) and one model seed.
% TODO: update after the scale-up (2,000 files, seeds 0-2).

The feature-based detector is a LightGBM GBDT on the EMBER v3 features with EMBER's baseline parameters (test ROC-AUC
0.9976). The raw-byte detector is the medium MalConv checkpoint released with EMBER-2024 (seed 0; test ROC-AUC 0.9913),
run on the same truncated files. Where a detector must decide, its threshold gives a 1\% false-positive rate on the
benign test files.

We explain the GBDT with three Shapley-value explainers: TreeSHAP in its path-dependent and interventional variants,
the latter with a background of 1{,}000 balanced training files, and KernelSHAP with 32{,}768 coalitions and a
balanced background of 16 files. We explain MalConv with Integrated Gradients over its byte embeddings (256 steps, the
padding embedding as baseline, summed over the embedding dimension).

\subsection{Measuring Faithfulness by Deletion}
\label{sec:XML:SetupDeletion}

An explanation is faithful if the inputs it ranks highest are the inputs the decision rests on. We test this as
explanations are usually tested~\cite{petsiuk2018rise}, by deleting the top-ranked inputs and measuring how much of
the detector's evidence disappears. Deleting a feature means setting it to a reference value
(Section~\ref{sec:XML:PushforwardFaithfulness}); deleting bytes means replacing them, and naive replacement breaks PE
files in ways that would make any comparison meaningless.

Our experiments use knowledge of the file format in exactly three places, each to protect an experiment's integrity
and never in the method itself: the deletion operator described here, the planted struct trigger of
Section~\ref{sec:XML:PullbackTrigger}, and the loader-safe modifications of the audit in Section~\ref{sec:XML:Audit}.
The deletion operator never changes the fields that define the file's layout (found by a sweep over 2{,}000 training
files) or the fields that break a particular file, replaces the remaining header fields with their most common value in
the training set, zeroes the content bytes, and verifies every perturbed file; none of the 23{,}000 perturbed files of
Section~\ref{sec:XML:PullbackFaithfulness} failed verification. Zeroing the header fields instead, or filling the
content with random bytes, serves as a robustness check.

A ranking is scored by deleting its top-ranked bytes at log-spaced budgets from 0.1\% to 50\% of the file,
re-extracting the features where the detector needs them, re-scoring the detector, and recording the \emph{evidence
removed}, how far the log-odds fall toward the other class, $s \cdot (f(x) - f(x'))$ with $s$ the sign of the original
log-odds. A ranking that finds the bytes the decision rests on removes much evidence already at small budgets. We
summarize a ranking by its evidence removed averaged over the budgets and then over the files, compare rankings file
by file with paired Wilcoxon tests, and count cost in feature extractions per file. Equal scores are common, among
bytes with identical provenance and among windows whose deletion changes nothing, so every ranking breaks ties by a
seeded random permutation rather than by position.

Any deletion produces files that a detector never saw in training, so a deletion curve measures how a model responds
to unusual inputs as well as what it relies on. We therefore compare rankings only against one another, at equal
budgets and under both fills, and rest no conclusion on deletion alone: the planted triggers
(Sections~\ref{sec:XML:PullbackTrigger} and~\ref{sec:XML:PushforwardTrigger}) give a ground truth without deleting
anything, and the audit (Section~\ref{sec:XML:Audit}) changes files only in ways the Windows loader accepts.

\subsection{Windows of Bytes}
\label{sec:XML:SetupWindows}

Several baselines explain the detectors over windows of bytes rather than single bytes, and comparing with them needs
$P$ over windows. Summing the rows of $P$ within each window $g$ of a partition $G$,
$P^G_{g,j} = \sum_{i \in g} P_{i,j}$, gives a provenance matrix over windows that satisfies the same axioms, and by
linearity, the pullback through $P^G$ is the pullback through $P$ summed within each window. Coarsening loses
separation: a window holds the bytes of many features, so its attribution can no longer be traced back to the features
that produced it, and the coarser the windows, the more features share each one. We use $P^G$ only to compare the
transfers with explainers that work over windows; it is not part of the method.
